% Generated draft: requires mathematical and source review.
\documentclass[11pt]{article}
\usepackage[margin=25mm]{geometry}
\usepackage{fontspec}
\setmainfont{Latin Modern Roman}
\usepackage{amsmath,amssymb,mathtools,unicode-math}
\setmathfont{Latin Modern Math}
\usepackage{enumitem,needspace,xurl,hyperref}
\setlength{\parindent}{0pt}
\setlength{\parskip}{5pt}
\setlength{\emergencystretch}{3em}
\newcommand{\N}{\mathbb{N}}
\newcommand{\Z}{\mathbb{Z}}
\newcommand{\Q}{\mathbb{Q}}
\newcommand{\R}{\mathbb{R}}
\newcommand{\C}{\mathbb{C}}
\newcommand{\F}{\mathbb{F}}
\newcommand{\A}{\mathbb{A}}
\newcommand{\PP}{\mathbb{P}}
\newcommand{\E}{\mathbb{E}}
\newcommand{\eps}{\varepsilon}
\newcommand{\dd}{\,\mathrm{d}}
\begin{document}
\section*{UnsolvedMath ID: 700675}
\textbf{Source problem code:} catalogue entry
\subsection{Definitions and mathematical statement}
Let \(B\) be an abelian group, written additively, and let \(A\leq\operatorname{Aut}(B)\). Under the usual fixed-point-free interpretation of a group of regular automorphisms, the hypothesis is
\[
\forall a\in A\setminus\{1\},\qquad \{b\in B:a(b)=b\}=\{0\}.
\]
The source does not define ``regular,'' so this interpretation requires confirmation. A group \(A\) is periodic if every element has finite order. For a prime \(p\), a periodic \(p\)-group is a group in which every element has order \(p^n\) for some integer \(n\geq0\). Its centre is
\[
Z(A)=\{a\in A:ax=xa\text{ for every }x\in A\}.
\]
The question is whether, for every abelian group \(B\) and every periodic subgroup \(A\leq\operatorname{Aut}(B)\) satisfying the displayed fixed-point condition, one necessarily has
\[
Z(A)\neq\{1\}.
\]
The parenthetical wording ``\(p\)-group'' also raises the corresponding question with \(A\) restricted to a periodic \(p\)-group, for an arbitrary prime \(p\). The record does not specify whether this is a separate restricted question or the intended domain of the whole question. It also does not explicitly require \(A\) or \(B\) to be nontrivial; the intended nontriviality hypotheses require clarification and have not been silently added.

The supplied editors' comment, dated 2001, reports an affirmative answer under the additional hypothesis that \(A\) contains an element of order \(2\) or \(3\), attributing this result to A. Kh. Zhurtov. This is a supplied historical qualification, not an independently verified theorem or a resolution of the unrestricted question.
\subsection{Short English statement}
The problem asks whether every periodic group of regular automorphisms of an abelian group has a nontrivial centre, and mentions a \(p\)-group version parenthetically. Interpreting regularity as the absence of nonzero fixed points for every nonidentity automorphism requires confirmation, as do the intended nontriviality hypotheses. A supplied 2001 comment reports an affirmative answer when the automorphism group contains an element of order \(2\) or \(3\).
\subsection{Sources}
\begin{itemize}
\item Local UnsolvedMath registry, permanent ID \(700675\), source problem code \texttt{catalogue entry}, title ``Kourovka Notebook Problem 10.60.'' The supplied registry labels the problem open; this label is not an independent verification of its present status.
\item A. I. Sozutov, problem 10.60, \emph{The Kourovka Notebook (Unsolved Problems in Group Theory)}, 10th issue (1986). The supplied background identifies the 21st issue edition (2026), arXiv:1401.0300v48. The collection description names E. I. Khukhro and V. D. Mazurov as editors. These bibliographic details are supplied attribution and have not been independently checked.
\item Editors' comment (2001), citing A. Kh. Zhurtov, \emph{Siberian Math. J.}, volume 41, no. 2 (2000), 268--275, for the reported affirmative answer when \(A\) contains an element of order \(2\) or \(3\). No article title is supplied.
\end{itemize}
\end{document}
